\subsection{Types of financial data}

\begin{qfdefinition}[Four types of financial data: fundamental, market, analytics and alternative data, with their timing, cost and uniqueness]\label{def:financial-data-types}
\qffield{Statement} López de Prado orders four ``essential types of financial data'' by increasing diversity, from fundamental to alternative data \citep[chap.~2, sec.~2.2, Table~2.1, pp.~23--24]{alpha-afml}:
\begin{itemize}\setlength{\itemsep}{0pt}
\item \emph{fundamental data}: assets, liabilities, sales, costs and earnings, macro variables; ``mostly accounting data, reported quarterly'' and ``reported with a lapse'' (p.~23);
\item \emph{market data}: prices, yields and implied volatilities, volume, dividends and coupons, open interest, quotes and cancellations, aggressor side; ``all trading activity that takes place in an exchange \dots\ or trading venue'' (p.~24);
\item \emph{analytics}: analyst recommendations, credit ratings, earnings expectations, news sentiment; ``derivative data, based on an original source'', processed for the user and ``not readily available from an original source'' (p.~25);
\item \emph{alternative data}: satellite and CCTV images, web searches, social media, metadata; ``primary information, that is, information that has not made it to the other sources'' (p.~25).
\end{itemize}
Following Kolanovic and Krishnamachari (2017), he divides alternative data into data from individuals, business processes and sensors (p.~25).

\qffield{Conditions} Practitioner descriptions as of 2018--2020; the 10TB figure undated and unsourced. \emph{Timing.} Fundamental data must be used only from the moment each point was released; López de Prado's example is Bloomberg, which indexes fundamental data by the last date in the report, ``which precedes the date of the release (often by 1.5 months)'' (p.~23). It is also backfilled (missing values assigned values ``unknown at that time'') and reinstated (corrections overwriting the first release), so a dataset that keeps only final values assigns information to dates when it was not known; some vendors keep every release, for example three values for a quarterly GDP figure (p.~24).

\emph{Cost and use.} Analytics may be costly, their methodology ``biased or opaque'', and ``you will not be the sole consumer''; alternative data raise cost and privacy concerns (p.~25). Velu, Hardy and Nehren add that the value of an alternative dataset depends on the user's resources and ``maybe more importantly---on the uniqueness of the data set. As more investors get access to it, the harder it becomes to extract meaningful alpha'' \citep[chap.~1, sec.~1.4.4, pp.~43--44]{s3-velu}.

\qffield{Practice} López de Prado's judgements are opinions, not findings: fundamental data being ``extremely regularized and low frequency'' leaves ``rather unlikely that there is much value left'', though it may help combined with other data (p.~24); FIX market data is ``appealing'' because it is ``not trivial to process'' and abundant, ``over 10TB being generated on a daily basis'' (p.~25, no source or market given); and ``a dataset may be useful if it annoys your data infrastructure team'', since competitors may not have tried, gave up or processed it wrongly (p.~25).

Velu, Hardy and Nehren list fundamental and analyst inputs with their own hedges (pp.~42--43): key ratios (EPS, P/E, P/B) for cross-sectional comparison; analyst recommendations, individually ``noisy'' but in aggregate a consensus whose changes ``can have a direct impact on price''; earnings estimates and surprises; holdings disclosed on Form~13F, where sudden changes ``might indicate changes in sentiment by sophisticated investor''; insider purchases and sales, ``thought to be an indicator of future stock price moves''; and credit rating downgrades, which ``generally have a negative impact on equity prices''. These are claims to test, not results.

\emph{Insider trades, tested}. Cohen, Malloy and Pomorski split insiders into ``routine'' traders, predictable from their own past trading pattern, and ``opportunistic'' ones, using Thomson Reuters insider filings for 1986--2007. Routine trades, more than half the universe, carry ``essentially zero'' abnormal returns. A portfolio following only opportunistic insiders earns value-weighted abnormal returns of 82 basis points a month \citep[working-paper version, abstract and introduction, pp.~5--6]{eq-cohen-malloy-pomorski-2012}. So Velu, Hardy and Nehren's hedge (``thought to be an indicator'') fits the whole universe, and the signal is concentrated in a subset identifiable in advance.

\emph{The numbers, with their uncertainty}. A routine trader is one who traded in the same calendar month for a set number of past years, classified at the start of each year from past history only; the opportunistic long--short portfolio (opportunistic buys minus sells) earns 82 basis points a month value-weighted ($t=2.15$) and 180 equal-weighted ($t=6.07$), against $-20$ ($t=-0.57$) and 43 ($t=1.73$) for the routine portfolio, on open-market trades from 1986 to 2007 \citep[NBER WP~16454, sec.~I, PDF p.~7]{eq-cohen-malloy-pomorski-2012}. The value-weighted figure, the investable one for large books, has a $t$ close to two; the equal-weighted figure is dominated by small firms. Its timing in a backtest still depends on when the filing became public (\hyperref[data:holdings-insider-filing-lags]{filing lags}). Thirteen-F holdings changes and rating downgrades are not tested in these sources.

\emph{Alternative data and forecast horizon}. Dessaint, Foucault and Fr\'esard argue that alternative data raise the benefit of short-term over long-term information. With sell-side analysts' forecasts and a new measure of informativeness by horizon, they confirm that alternative data make ``short-term forecasts more informative but long-term forecasts less so'' \citep[working-paper version of 20 July 2021, not the 2024 article, abstract and sec.~I]{rf-dessaint-foucault-fresard-2024}. Farboodi and Veldkamp's growth model reaches a related theoretical point: better data technology shifts traders toward extracting information from order flow, which ``does crowd out production of fundamental information'', yet does not lower price informativeness \citep[NBER WP 23457, May 2017, not the 2020 article, abstract]{rf-farboodi-veldkamp-2020}. Whether a given dataset's predictive value falls as access spreads (Katona et al.) is a separate question.

\qffield{Pitfalls} The Bloomberg indexing lag is specific to that example, not a universal delay for fundamental data. The two source taxonomies differ; holdings and insider-data features must account for filing lags.

\qffield{See also} \hyperref[data:trading-reference-data]{trading reference data}; \hyperref[data:market-data-levels]{market data levels}; \hyperref[data:derived-market-statistics]{derived market statistics}; \hyperref[meth:activity-sampled-bars]{activity-sampled bars}.

\qffield{Further reading} \nolinkurl{wide-kolanovic-big-data} and \nolinkurl{rf-monk-prins-rook-2019}, and the evidence on how alternative-data edges decay.

\end{qfdefinition}
